% vim: set textwidth=78 autoindent:

\section{Usar complementos externos de Python para QGIS}\label{sec:external_plugins}\index{plugins}

% when the revision of a section has been finalized, 
% comment out the following line:
% \updatedisclaimer

Los complementos externos de QGIS están escritos en Python. Se guardan en los repositorios de QGIS
'Oficial' o 'Contribuido por usuarios' o bien en varios otros repositorios externos mantenidos por los autores individuales.
La tabla \ref{tab:external_plugins} muestra una lista de los complementos 'Oficiales' disponibles actualmente, con una breve descripción.
La documentación detallada sobre el uso, versión mínima de QGIS, página web, autores y otra información importante se
proporciona con los propios complementos y no está incluida en este manual.
\footnote{Las actualizaciones de complementos del núcleo pueden estar disponibles en este repositorio como capas externas.} 
\footnote{Los complementos fTools, exportación a Mapserver y el Instalador de complementos son complementos de Python, 
pero también son parte de las fuentes de QGIS y se cargan y activan automáticamente dentro del Administrador de complementos 
de QGIS (ver Sección~\ref{sec:load_external_plugin}).}

Encontrará una lista actualizada de complementos 'Oficiales' en el Repositorio Oficial de QGIS en
\url{http://qgis.osgeo.org/download/plugins.html}. Esta lista también está disponible automáticamente desde el 
\filename{Instalador de complementos} via \dropmenuopttwo{plugin_installer}{Obtener complementos de Python...}.

\begin{table}[H]
\centering
\caption{Complementos externos moderados actuales de QGIS Plugins}\label{tab:external_plugins}\medskip
\small
 \begin{tabular}{|l|l|p{4in}|}
\hline \textbf{Icono} & \textbf{Complemento externo} & \textbf{Descripción}\\
\hline
\includegraphics[width=0.7cm]{zoom2point_icon}
 & Zum a punto \index{plugins!Zoom To Point} & Acerca la imagen a las coordenadas especificadas en el diálogo de 
entrada. Puede especificar el nivel de zum así como controlar la extensión de la vista.\\
\hline
\includegraphics[width=0.7cm]{plugin_installer}
 & Instalador de complementos \index{plugins!Plugin Installer} & El instalador de complementos de Python más reciente.\\
\hline

\end{tabular}
\end{table}

Se puede encontrar una descripción detallada del procedimiento de instalación para complementos externos de python en
la Sección \ref{sec:load_external_plugin}.

\begin{Tip} \caption{\textsc{Añadir más repositorios}}
\qgistip{Para añadir el repositorio 'Contribuido por usuarios' y/o varios repositorios de autores externos, abra el
Instalador de complementos (\mainmenuopt{Complementos} > \dropmenuopttwo{plugin_installer}{Obtener complementos de Python...}),
vaya a la pestaña \tab{Repositorios} y pulse \button{Añadir repositorios de terceros}. 
Si no quiere uno o más de los 
repositorios añadidos, los puede deshabilitar mediante el botón 
\button{Editar...} o eliminarlos completamente con el botón \button{Borrar}.
}
\end{Tip}

